% ============================================================
\chapter{Vorgehen, Zeitplan und Projektorganisation}
% ============================================================

Die Leistungserbringung erfolgt in mehreren aufeinander aufbauenden
Projektphasen. Der konkrete Phasen- und Meilensteinplan wird nach Auswahl
der Leistungsvariante gemeinsam mit der Universitätsstadt Tübingen
finalisiert und im Projekthandbuch dokumentiert.

Das Projekt orientiert sich am V-Modell~XT in der Projekttypvariante
\enquote{AN-Projekt mit Weiterentwicklung} und verbindet dessen strukturierte
Phasen-, Rollen- und Qualitätssicherungslogik mit einer iterativen Umsetzung
in regelmäßigen Entwicklungs- und Abstimmungsschleifen.

\section{Projektphasen und Meilensteine}

\paragraph{P1 -- Grundlagen und Konzeption (ca. Monat 1)}
Aktualisierung der Systemarchitektur, Erweiterung der Entscheidungslogik und
des Nutzerflusses, Berücksichtigung relevanter Personengruppen und
Haushaltskonstellationen sowie Erarbeitung der Datenschutz-, Compliance-
und UX-Grundlagen.\\
\textbf{Meilenstein 1: Konzeptionsgrundlage abgestimmt}

\paragraph{P2 -- Review und Integration (ca. Monat 2)}
Review von Systemarchitektur und Entscheidungslogik, Integration zentraler
Logikbestandteile in KIWo, UX-Update sowie frühe Qualitätssicherung durch
ein Schwachstellen-Screening.\\
\textbf{Meilenstein 2: KIWo grundsätzlich nutzbar}

\paragraph{P3 -- Begleitete Tests und Weiterentwicklung (ca. Monat 3 bis 4)}
Begleitete Nutzertests, Feedback aus Verwaltung und Beratungsstellen,
Auswertung der Rückmeldungen, fachliche Nachschärfungen, technische
Weiterentwicklung und Sicherheitsprüfung vor breiterer Nutzung.\\
\textbf{Meilenstein 3: Begleitete Tests ausgewertet und System weiterentwickelt}\\
\textbf{Meilenstein 4: Sicherheitsprüfung vor breiterer Nutzung abgeschlossen}

\paragraph{P4 -- Stabilisierung und Onboarding (ca. Monat 5)}
Nacharbeiten aus den begleiteten Tests, Bug-Fixes, Stabilisierung der
Systemversion und Vorbereitung der Testgruppen. In dieser Phase werden
grundsätzlich keine neuen Kernfunktionen aufgenommen.\\
\textbf{Meilenstein 5: Breiterer Testbetrieb vorbereitet}

\paragraph{P5 -- Nutzertests, Evaluation und Abschluss (ca. Monat 6 bis 7)}
Durchführung des Testbetriebs, Bearbeitung akuter technischer Probleme,
Sammlung und Auswertung von Evaluationsdaten sowie Erstellung der
Abschlussdokumentation.\\
\textbf{Meilenstein 6: Evaluation und Abschlussbericht vorgelegt}

\section{Auszahlungskonzept}

Die Auszahlung erfolgt meilensteinbezogen. Die genaue Ausgestaltung wird
nach Auswahl der Leistungsvariante gemeinsam festgelegt. Die folgende
Struktur dient als Vorschlag für eine nachvollziehbare Auszahlung entlang
prüfbarer Projektstände.

\begin{tabularx}{\textwidth}{@{} X X r @{}}
  \toprule
  \tblhead{Zahlungszeitpunkt} & \tblhead{Nachweisbarer Auslöser} & \tblhead{Anteil} \\
  \midrule
  Projektstart
    & Beauftragung erfolgt; Projekt ist organisatorisch gestartet und
      die Konzeptionsphase beginnt.
    & 25~\% \\[0.4em]
  Konzeptionsgrundlage abgestimmt
    & Systemarchitektur, Entscheidungslogik, Nutzerfluss,
      Datenschutz-/Compliance-Grundlagen und UX-Grundlage liegen als
      abgestimmter Arbeitsstand vor.
    & 20~\% \\[0.4em]
  KIWo grundsätzlich nutzbar
    & Zentrale Logikbestandteile sind integriert; ein lauffähiger
      KIWo-Stand kann demonstriert und für begleitete Tests genutzt werden.
    & 25~\% \\[0.4em]
  Breiterer Testbetrieb vorbereitet
    & Begleitete Tests wurden ausgewertet, wesentliche Nacharbeiten
      umgesetzt, Sicherheitsprüfung abgeschlossen und
      Onboarding-Unterlagen vorbereitet.
    & 20~\% \\[0.4em]
  Projekt abgeschlossen
    & Testbetrieb, Evaluation und Abschlussdokumentation sind
      abgeschlossen; Ergebnisse sind für die Bewertung von Nutzen,
      Machbarkeit und Skalierung aufbereitet.
    & 10~\% \\
  \midrule
  \textbf{Gesamt} & & \textbf{100~\%} \\
  \bottomrule
\end{tabularx}

\section{Qualitätssicherung und Projektorganisation}

Die Qualitätssicherung im Projekt KIWo~2.0 verfolgt das Ziel, die fachliche
Korrektheit, technische Stabilität, Sicherheit und Nutzbarkeit des Systems
über den gesamten Projektverlauf hinweg systematisch abzusichern. Sie ist
als kontinuierlicher, projektbegleitender Prozess angelegt und orientiert
sich an den Vorgaben des V-Modell~XT.

Auf technischer Ebene umfasst die Qualitätssicherung insbesondere
systematische Code-Reviews, automatisierte Tests sowie ein fortlaufendes
Schwachstellen-Screening. Ergänzend wird vor der breiteren Nutzung ein
externer Penetrationstest durchgeführt, einschließlich eines Re-Tests zur
Verifikation der behobenen Findings. Damit wird sichergestellt, dass
identifizierte Schwachstellen nicht nur dokumentiert, sondern auch
nachweislich geschlossen werden.

Auf fachlicher und nutzerorientierter Ebene wird die Qualitätssicherung
durch die Begleitung der Test- und Pilotbetriebsphase ergänzt. Hierzu
gehören strukturierte Nutzertests, die Einbindung von Antragstellenden,
Mitarbeitenden der Verwaltung sowie Beratungsstellen und die systematische
Erhebung und Auswertung von Praxisrückmeldungen.

Risiken werden im Rahmen des Projektmanagements fortlaufend identifiziert,
bewertet und gemeinsam mit der Universitätsstadt Tübingen abgestimmt. Eine
ausführliche und variantenspezifische Ausarbeitung der Qualitätssicherungs-
maßnahmen sowie des Risikomanagements erfolgt im Rahmen der
Projektinitialisierungsphase und wird Bestandteil des Projekthandbuchs.

Die Projektleitung ist in Tübingen ansässig und steht der Universitätsstadt
als örtliche Schnittstelle zur Verfügung. Systemarchitektur, Qualitäts-
sicherung und Entwicklung sind in Köln kolokiert. Externe Spezialisten
werden in den Bereichen UX-Design, Datenschutz, IT-Sicherheit und KI-Recht
eingebunden.

\section{Mitwirkung der Universitätsstadt Tübingen}

Die Durchführung des Folgeprojekts setzt eine enge, pragmatische
Zusammenarbeit mit der Universitätsstadt Tübingen voraus. Die Mitwirkung
orientiert sich an der bewährten Zusammenarbeit aus der bisherigen
Projektphase und wird nach Auswahl der Leistungsvariante konkret abgestimmt.

Zentral ist insbesondere die fachliche Einbindung der Wohngeldstelle.
Diese unterstützt nach Möglichkeit bei der Erweiterung der
Entscheidungslogik auf weitere Personengruppen und Haushaltskonstellationen,
bei fachlichen Reviews der umgesetzten Logik sowie bei der Bewertung
praxisnaher Testfälle.

Darüber hinaus unterstützt die Stadt das Projekt insbesondere durch:

\begin{itemize}
  \item Vermittlung zu Beratungsstellen für den Testbetrieb.
  \item Nutzung geeigneter Informationskanäle zur Ansprache potenzieller Testnutzer:innen.
  \item Teilnahme an projektbezogenen Abstimmungen und Lenkungsentscheidungen.
  \item Notwendige organisatorische und rechtliche Abstimmungen, etwa zu Datenschutz, Auftragsverarbeitung und Testbetrieb.
\end{itemize}
